%% ma: L10-B14-0007
%% lop: 10
%% chuong: 5
%% bai: 14
%% dang: 10.14.1
%% loai: TL
%% muc: NB
%% dap_an: "-"
%% nguon: "chuyen-de-cac-so-dac-trung-cua-mau-so-lieu-khong-ghep-nhom-toan-10-chuong-trinh-moi.pdf (GV. Phan Nhật Linh) (P04), Bài 03 Dạng 1 Bài tập tự luận 7"
%% kien_thuc: "Lớp 10 Bài 14 (khoảng biến thiên)"
%% trang_thai: chinh-thuc
%% ghi_chu: "Lời giải gốc ghi nhầm 'số lớn nhất là 9 và bé nhất là 5' nhưng kết quả 3,5 đúng. Tự duyệt theo yêu cầu giáo viên 22/10/2026."
\begin{ex}
Điểm học kì I của một học sinh được cho bởi bảng số liệu sau:
\begin{center}
\begin{tabular}{|c|c|c|c|c|c|c|c|}\hline
$6$ & $6$ & $6$ & $7$ & $7$ & $8$ & $8{,}5$ & $9{,}5$\\\hline
\end{tabular}
\end{center}
Tìm khoảng biến thiên của mẫu số liệu.
\loigiai{Giá trị lớn nhất là $9{,}5$, nhỏ nhất là $6$ nên $R=9{,}5-6=3{,}5$.}
\end{ex}
